\chapter{The \solver{} Solver}
\label{chap:method}

\section{Design overview}
\label{sec:method-overview}

\solver{} is a block structured adaptive lattice kinetic solver written for
distributed memory machines. The domain is partitioned into blocks of
$32^{3}$ lattice sites, each block sits at a refinement level $\ell$ with
spacing $\Delta x_\ell = \Delta x_0 2^{-\ell}$, and the hierarchy is limited to
four levels in the production configuration. Figure~\ref{fig:pipeline} shows the
arrangement of a time step.

\begin{figure}[t]
\centering
\begin{tikzpicture}[
  node distance=8mm,
  stage/.style={draw,rounded corners=2pt,minimum width=27mm,minimum height=9mm,
                align=center,font=\small},
  arrow/.style={-{Stealth[length=2mm]},thick}
]
  \node[stage] (collide) {Collide\\ \footnotesize moment space};
  \node[stage,right=of collide] (stream) {Stream\\ \footnotesize shift by $\mathbf{c}_i$};
  \node[stage,right=of stream] (bc) {Boundaries\\ \footnotesize walls and periodic};
  \node[stage,below=12mm of bc] (coarse) {Restrict\\ \footnotesize fine to coarse};
  \node[stage,left=of coarse] (fine) {Prolong\\ \footnotesize coarse to fine};
  \node[stage,left=of fine] (regrid) {Regrid\\ \footnotesize every $n_r$ steps};

  \draw[arrow] (collide) -- (stream);
  \draw[arrow] (stream) -- (bc);
  \draw[arrow] (bc) -- (coarse);
  \draw[arrow] (coarse) -- (fine);
  \draw[arrow] (fine) -- (regrid);
  \draw[arrow] (regrid) -- ++(0,12mm) -| (collide);

  \begin{scope}[on background layer]
    \node[draw,dashed,inner sep=3mm,fit=(coarse)(fine)] (couple) {};
  \end{scope}
  \node[font=\footnotesize,above=1mm of couple.north east,anchor=south east]
    {interface coupling, Section~\ref{sec:method-coupling}};
\end{tikzpicture}
\caption[Structure of a \solver{} time step]{Structure of a \solver{} time step.
Collision and streaming are level local. The interface coupling runs after
boundary treatment and before the next collision, and the regridding pass runs
once every $n_r$ steps.}
\label{fig:pipeline}
\end{figure}

Levels are advanced with subcycling in time, so a block at level $\ell$ takes
$2^{\ell}$ steps for every step taken at level zero. This keeps the lattice
Mach number and the relaxation rates constant across the hierarchy, which is the
condition under which the interface rescaling of
Section~\ref{sec:method-coupling} is exact \\citep{ternes2021subcycling}.

\section{Refinement criterion}
\label{sec:method-criterion}

A conventional criterion refines where the velocity gradient magnitude, or the
vorticity, exceeds a threshold. In a stratified layer this is the wrong
diagnostic. The strongest shear sits in the layers where stratification has
already suppressed overturning, and refining those layers buys resolution in a
region that carries almost no scalar flux while starving the intermittent sheets
that carry nearly all of it \citep{vestergaard2022ozmidov}.

\solver{} keys instead on the resolved fraction of the Ozmidov scale. For each
block $B$ the solver forms the block dissipation and buoyancy frequency,
\begin{equation}
  \varepsilon_B = \frac{2\nu}{|B|} \sum_{\mathbf{x} \in B} S_{\alpha\beta} S_{\alpha\beta},
  \qquad
  N_B^{2} = \beta g \left\langle \partial_z \theta + \Gamma \right\rangle_B,
  \label{eq:block-diagnostics}
\end{equation}
where $S_{\alpha\beta}$ is recovered directly from the non equilibrium stress
of~\eqref{eq:stress} without a finite difference, and then the block refinement
indicator
\begin{equation}
  \chi_B = \frac{1}{\Delta x_\ell}
    \left( \frac{\varepsilon_B}{\max(N_B^{2}, N_{\min}^{2})^{3/2}} \right)^{1/2}.
  \label{eq:indicator}
\end{equation}
A block refines when $\chi_B < \chi_{\mathrm{ref}}$ and coarsens when
$\chi_B > \chi_{\mathrm{coar}}$, with the production thresholds
$\chi_{\mathrm{ref}} = 2$ and $\chi_{\mathrm{coar}} = 8$ chosen from the
sensitivity study of Section~\ref{sec:val-thresholds}. The floor
$N_{\min}^{2} = 10^{-4} N_0^{2}$ prevents the indicator from diverging in
neutral pockets, where the criterion falls back to a dissipation threshold.

The two thresholds are separated by a factor of four to prevent a block from
oscillating between levels, which in an early version of the solver consumed
eleven percent of the total runtime in regridding alone.

\begin{algorithm}[t]
\SetAlgoLined
\KwIn{Block hierarchy $\mathcal{H}$, thresholds $\chi_{\mathrm{ref}}, \chi_{\mathrm{coar}}$, maximum level $\ell_{\max}$}
\KwOut{Updated hierarchy $\mathcal{H}'$}
\ForEach{block $B \in \mathcal{H}$}{
  compute $\varepsilon_B$ and $N_B^{2}$ from~\eqref{eq:block-diagnostics}\;
  compute $\chi_B$ from~\eqref{eq:indicator}\;
  \uIf{$\chi_B < \chi_{\mathrm{ref}}$ \textbf{and} $\ell(B) < \ell_{\max}$}{
    mark $B$ for refinement\;
  }
  \uElseIf{$\chi_B > \chi_{\mathrm{coar}}$ \textbf{and} $\ell(B) > 0$}{
    mark $B$ for coarsening\;
  }
}
propagate refinement marks to face neighbours to enforce a two to one level
ratio\;
withdraw coarsening marks on any block with a refined neighbour\;
\ForEach{marked block $B$}{
  \uIf{$B$ refines}{
    allocate eight children, prolong $f_i$ and $g_j$ by
    Equation~\eqref{eq:prolong}\;
  }
  \Else{
    restrict children by Equation~\eqref{eq:restrict}, release them\;
  }
}
rebalance blocks across ranks by space filling curve order\;
\caption{Regridding pass, executed every $n_r = 200$ coarse steps.}
\label{alg:regrid}
\end{algorithm}

\section{Interface coupling}
\label{sec:method-coupling}

The coupling problem is that a population is not a field. Splitting
$f_i = f_i^{\mathrm{eq}} + f_i^{\mathrm{neq}}$, the equilibrium part is a
function of $\rho$ and $\mathbf{u}$ and so transfers between levels by ordinary
interpolation. The non equilibrium part carries the stress
of~\eqref{eq:stress}, which depends on the relaxation rate, and the relaxation
rate must change with the grid spacing if the physical viscosity is to stay
fixed.

Requiring that the stress $\Pi_{\alpha\beta}^{(1)}$ be continuous across the
interface, and using~\eqref{eq:viscosity} at both levels, gives the rescaling
factor
\begin{equation}
  f_{i,\ell+1}^{\mathrm{neq}} = \frac{1}{2}
    \frac{s_{\nu,\ell}}{s_{\nu,\ell+1}} \, f_{i,\ell}^{\mathrm{neq}},
  \qquad
  \frac{s_{\nu,\ell+1}}{s_{\nu,\ell}}
    = \frac{2}{1 + \tfrac{1}{2} s_{\nu,\ell}}.
  \label{eq:rescale}
\end{equation}
Prolongation to a child block is then
\begin{equation}
  f_{i,\ell+1}(\mathbf{x}) = f_{i,\ell+1}^{\mathrm{eq}}\!\left[ \mathcal{I}\rho, \mathcal{I}\mathbf{u} \right](\mathbf{x})
    + \frac{1}{2}\frac{s_{\nu,\ell}}{s_{\nu,\ell+1}} \, \mathcal{I} f_{i,\ell}^{\mathrm{neq}}(\mathbf{x}),
  \label{eq:prolong}
\end{equation}
with $\mathcal{I}$ a tricubic interpolation operator, and restriction is the
volume average with the inverse factor,
\begin{equation}
  f_{i,\ell}(\mathbf{x}) = f_{i,\ell}^{\mathrm{eq}}\!\left[ \langle \rho \rangle, \langle \mathbf{u} \rangle \right](\mathbf{x})
    + 2 \frac{s_{\nu,\ell+1}}{s_{\nu,\ell}} \left\langle f_{i,\ell+1}^{\mathrm{neq}} \right\rangle (\mathbf{x}).
  \label{eq:restrict}
\end{equation}

Because $\mathcal{I}$ is linear and reproduces constants, and because the
equilibrium is constructed from interpolated moments rather than interpolated
populations, mass and momentum are conserved by construction and the second
order moment is conserved by~\eqref{eq:rescale}. The scalar population uses the
same construction with $s_\nu$ replaced by $1/\tau_g$. Section~\ref{sec:val-interface}
measures the consequence of omitting the rescaling factor, which is the
configuration used by most earlier refined lattice solvers.

\section{Parallel decomposition and cost}
\label{sec:method-parallel}

Blocks are ordered along a Hilbert curve \\citep{kalliokoski2020hilbert} and
distributed to ranks in contiguous runs weighted by $8^{\ell}$, which accounts for both the site count and the
subcycling factor. Halo exchange is a single non blocking neighbour exchange per
level per substep. Table~\ref{tab:cost-breakdown} gives the measured cost
distribution for the reference production case on 512 nodes.

\begin{table}[t]
\centering
\caption[Cost breakdown for a production case]{Measured cost breakdown per
coarse time step for case S07 on 512 nodes of the Noorderlicht system. Times are
medians over 2000 steps. The regridding column is amortised over the $n_r = 200$
step interval.}
\label{tab:cost-breakdown}
\begin{tabular}{lrrr}
\toprule
Stage & Time (ms) & Share (\%) & Bytes moved (MB) \\
\midrule
Collision, momentum   & 41.7 & 46.2 & 0 \\
Collision, scalar     &  8.3 &  9.2 & 0 \\
Streaming             & 14.9 & 16.5 & 0 \\
Halo exchange         & 12.1 & 13.4 & 184 \\
Interface coupling    &  9.6 & 10.6 & 27 \\
Boundary conditions   &  1.8 &  2.0 & 0 \\
Regridding, amortised &  1.9 &  2.1 & 41 \\
\midrule
Total                 & 90.3 & 100.0 & 252 \\
\bottomrule
\end{tabular}
\end{table}

The momentum collision dominates, as it should for a compute bound kinetic
scheme, and the interface coupling costs about a tenth of the step. That is the
price of adaptivity, and Section~\ref{sec:results-cost} shows it against the
factor of twenty three reduction in site count that adaptivity buys.
